% ============================================================
% Capítulo 1 — Introdução
% ============================================================

\chapter{Introdução}\label{cap:introducao}

\section{Contextualização}\label{sec:contextualizacao}

A internet acelerou a circulação de informação a um ponto sem precedentes e,
com ela, a de desinformação% TODO: inserir referência aqui
. Notícias falsas já existiam em jornais, rádio e televisão, mas as redes
sociais transformaram sua escala e velocidade. Conteúdos falsos são criados e
compartilhados por razões que vão do lucro financeiro à manipulação política,
passando pela exploração da economia da atenção. O viés de confirmação funciona
como acelerador: leitores tendem a compartilhar aquilo que reforça suas crenças
sem verificar a origem. As consequências são concretas e documentadas: protestos,
violência física, linchamentos, crimes de ódio, perseguição de minorias, erosão
da confiança em instituições públicas, pânico financeiro e boicotes comerciais
sem fundamento.

A pandemia de COVID-19 ilustra bem esse fenômeno. No Brasil, supostos
``remédios milagrosos'' sem qualquer respaldo científico foram promovidos em
larga escala. Nos Estados Unidos, a desinformação gerou danos físicos diretos:
segundo relatório do \textit{Centers for Disease Control and Prevention}
\cite{cdc2020}, 39\% dos adultos entrevistados relataram ter utilizado
desinfetantes e alvejantes de forma inadequada, motivados por falsas narrativas
de prevenção contra o vírus. Esses episódios evidenciam que a desinformação se
propaga mais rápido do que a capacidade humana de verificá-la, tornando
necessárias soluções automatizadas de detecção.

O padrão é estrutural. \citeonline{vosoughi2018} demonstraram que notícias
falsas se propagam de forma significativamente mais ampla, rápida e profunda do
que notícias verdadeiras nas redes sociais. Os algoritmos das plataformas agravam
esse cenário ao priorizar conteúdo com alta taxa de interação, amplificando
publicações que exploram reações emocionais independentemente de sua veracidade.
O problema da desinformação, portanto, não se restringe ao conteúdo, reside
também na topologia de propagação: quem compartilha o quê, para quem e com que
velocidade. Compreender essa estrutura é pré-requisito para uma abordagem eficaz
de detecção automatizada.


\section{Motivação}\label{sec:motivacao}

A verificação manual de veracidade não escala. O volume de postagens em redes
sociais opera em escala de \textit{Big Data}, com fluxos massivos gerados em
tempo real. Checar fontes e cruzar informações demanda um tempo que a velocidade
de propagação não concede. Produzir uma mentira é instantâneo; desmenti-la é
lento e custoso. Essa assimetria inviabiliza a curadoria humana e impõe a adoção
de modelos computacionais de detecção.

As abordagens tradicionais baseiam-se em Processamento de Linguagem Natural (PLN)
e extraem padrões semânticos e sintáticos associados à falsidade a partir do
conteúdo das publicações. Porém, essas técnicas falham diante de textos falsos
com alta sofisticação linguística, que passam despercebidos por classificadores
textuais. O que tais modelos desconsideram é a dinâmica de interação da rede: os
perfis que compartilham a informação, suas conexões e a velocidade de
disseminação. A topologia de propagação, conforme consolidado na literatura,
carrega sinal preditivo equivalente ou superior ao do próprio conteúdo textual
\cite{vosoughi2018,bian2020}. Redes Neurais em Grafos (\textit{Graph Neural
Networks} --- GNNs) exploram justamente esse sinal: esses modelos aprendem de
forma conjunta os atributos textuais, associados aos vértices do grafo, e a
topologia de propagação, representada pelas arestas, superando as limitações de
abordagens exclusivamente textuais.

A aplicação de GNNs na detecção de notícias falsas já foi investigada em
plataformas como Twitter e Facebook \cite{monti2019,dou2021}, mas o Bluesky
permanece praticamente inexplorado. Aberta ao público em 2023 e fundamentada na
arquitetura do AT Protocol (\textit{Authenticated Transfer Protocol}), a
plataforma cresceu rapidamente em decorrência de migrações em massa provocadas
por mudanças nas políticas de moderação de concorrentes. Sua infraestrutura
descentralizada e a transparência na disponibilização de dados criam um ambiente
propício para pesquisa acadêmica, ainda que faltem estudos que apliquem
modelagem baseada em grafos à sua topologia de propagação. Este trabalho preenche
essa lacuna ao integrar representações vetoriais textuais (\textit{embeddings}
extraídos via BERT) a Redes Convolucionais em Grafos (GCNs) para classificar
postagens no Bluesky. A contribuição é dupla: avançar nos métodos de detecção
automatizada de desinformação e ampliar a compreensão sobre como conteúdo falso
se propaga em plataformas descentralizadas emergentes.
